\chapter{A Map of All Markets}\label{ch:m3:a-map-of-all-markets}

The foreign exchange market trades close to ten trillion dollars a day; the most active bitcoin perpetual on
the largest crypto venue, under twenty billion; the weather contracts of \cref{ch:m3:freight-weather-and-other-corners}, far less. The first
three books of this series have described dozens of markets one at a time, each with its own instruments,
institutions, hours and habits. This last chapter puts them on one page: how big they are, how they are
organised, when they trade, who trades them, and what kinds of strategies they reward. It is a map for a
trading firm deciding where to go next, and for a reader deciding which of the later books to open.

\section{Size}

\begin{definition}[Turnover velocity]\label{def:m3:a-map-of-all-markets:velocity}
The \emph{turnover velocity}\index{turnover velocity} of a market is its daily turnover divided by the amount of the asset
outstanding (or, for derivatives, the open interest): the share of the stock that changes hands each day.
\end{definition}

Size alone misleads: a market can be huge and sleepy or small and frantic. The US Treasury market trades about USD~1.2 trillion
a day against USD~31.8 trillion outstanding, a velocity of 3.8\% a day. The largest bitcoin perpetual on one venue traded about
USD~17.6 billion in the 24 hours to the afternoon of 24 September 2026 against open interest of about 96{,}000 bitcoin, some USD~8
billion: a velocity above two, so that its open interest turned over twice a day. Velocity says what kind of trading a market rewards:
low velocity, the patient provision of balance sheet (One Quant Book 2); high velocity, speed and inventory turnover.

\begin{dated}{2026-09}{Market sizes on one scale}\label{dat:m3:a-map-of-all-markets:sizes}
The BIS Triennial Survey put average daily FX turnover at USD~9.6 trillion in April 2025. SIFMA reported US fixed-income trading of USD~1{,}659.2
billion a day, of which US Treasuries USD~1{,}203.6 billion, for 2026 to August, with USD~31.8 trillion of Treasuries outstanding. The FIA counted
119.29 billion exchange-traded futures and options contracts worldwide in 2025. On 24 September 2026 Binance's public endpoints reported 24-hour
quote volume of USD~17.65 billion in its BTCUSDT perpetual and USD~1.90 billion in BTCUSDT spot, and open interest of 96{,}463 bitcoin in the
perpetual; Deribit's bitcoin perpetual reported USD~1.58 billion of 24-hour volume. In 2025 the LME's futures and options averaged 717{,}334 lots a
day, and EPEX SPOT's power markets traded 917.5~TWh over the year.
\end{dated}

\begin{omfigure}
\begin{tikzpicture}
\begin{axis}[width=11cm,height=5.2cm,xbar,bar width=9pt,xmode=log,log basis x=10,
  xlabel={daily turnover, USD billion (log scale)},
  tick label style={font=\small},label style={font=\small},
  ymin=0.4,ymax=6.6,ytick={1,2,3,4,5,6},
  yticklabels={FX (Apr 2025),US fixed income,US Treasuries,BTC perp (Binance),BTC spot (Binance),BTC perp (Deribit)},
  y dir=reverse,xmin=1,xmax=20000,grid=major,grid style={gray!25},
  table/col sep=comma]
\addplot[fill=omThm!60,draw=omThm] table[x=usd_bn,y=idx]{figdata/markets-3/29-a-map-of-all-markets/sizes.csv};
\end{axis}
\end{tikzpicture}

\omcaption{Daily turnover of six markets on one logarithmic scale. The measures differ (the BIS counts all FX instruments net of double counting
in April 2025; SIFMA averages US trading for 2026 to August; the crypto figures are one day's volume on single venues), and the gap between the
largest and the smallest is nearly four orders of magnitude. Data: \cref{dat:m3:a-map-of-all-markets:sizes}.}\label{fig:m3:a-map-of-all-markets:sizes}
\end{omfigure}

\section{Structure}

Every market in these books sits somewhere between two poles. At one, a central limit order book on an exchange, open to many, with a central
counterparty behind every trade: equity and futures markets (One Quant Book 1), power day-ahead auctions and crypto venues in this book. At the other,
dealers quoting to clients by request for quote, bilateral credit and settlement: most bonds, swaps, FX forwards and physical commodities (One Quant
Book 2 and the first part of this book). Between them sit the hybrids: all-to-all bond platforms, cleared swaps traded on electronic platforms,
brokered commodity markets, and the \omterm{def:m3:automated-market-makers:amm}{automated market makers} and \omterm{def:m3:on-chain-order-books-and-perpetual-dexs:perpdex}{on-chain order books} of the crypto chapters, which remove the dealer and the
clearing house at once and put a protocol in their place.

The position of a market between the poles decides who can trade it and how. A firm with technology and little balance sheet goes to the order books;
a firm with balance sheet and client relationships goes to the dealer markets; a firm that can read a protocol goes to the chains. The trend of the last
twenty years has been towards the order-book pole, electronic and cleared, but the dealer markets remain among the largest by value, FX first.

\section{Hours}

\begin{definition}[Follow-the-sun trading]\label{def:m3:a-map-of-all-markets:sun}
\emph{Follow-the-sun trading}\index{follow-the-sun trading} is the organisation of a book so that it is managed around the clock by desks in successive time
zones, each handing its positions, orders and open issues to the next as its day ends.
\end{definition}

Markets have hours, and the hours have fixed moments at which prices are set and positions are marked. \Cref{fig:m3:a-map-of-all-markets:day} places some of
the fixed moments met in this book on one UTC day in the northern summer: the eight-hourly funding of perpetuals (\cref{ch:m3:perpetual-futures}), the 08:00
UTC expiry of crypto options (\cref{ch:m3:dated-futures-options-and-etfs}), the European day-ahead auction at 12:00 Central European Time
(\cref{ch:m3:power-markets-design}), the publication of the LME's official prices from the second Ring session (\cref{ch:m3:metals}), the US equity core session
from 9:30 to 16:00 New York time, and the hour from 15:00 to 16:00 New York time from which the bitcoin reference rate for US funds is computed. Around
these moments, electronic futures and FX run almost around the clock on weekdays: CME Group's FX market on Globex, for instance, trades from
Sunday 17:00 to Friday 16:00 Chicago time with a one-hour break each day from 16:00.

\begin{omfigure}
\centering
\begin{tikzpicture}[font=\footnotesize,x=0.47cm,y=0.62cm]
\draw[thick,->] (0,0) -- (24.6,0) node[right] {UTC};
\foreach \h in {0,4,8,12,16,20,24} {\draw (\h,0.12) -- (\h,-0.12) node[below] {\h:00};}
\fill[omThm!25] (0,0.3) rectangle (24,0.7); \node[anchor=west] at (0.1,0.5) {crypto: continuous, every day};
\foreach \h in {0,8,16} {\draw[omThm,very thick] (\h,0.2) -- (\h,0.8);}
\node[anchor=west] at (16.1,0.95) {perpetual funding at 00, 08, 16};
\fill[omExo!40] (8,1.3) rectangle (8.15,1.7); \node[anchor=east] at (7.9,1.5) {crypto options expiry 08:00};
\fill[omDef!40] (10,2.0) rectangle (10.15,2.4); \node[anchor=east] at (9.9,2.2) {EU day-ahead auction 10:00};
\fill[omProp!40] (11.33,2.7) rectangle (12.42,3.1); \node[anchor=east] at (11.2,2.9) {LME official prices 11:20--12:25};
\fill[omThm!40] (13.5,3.4) rectangle (20,3.8); \node[anchor=east] at (13.4,3.6) {US equities core session 13:30--20:00};
\fill[omExo!60] (19,4.1) rectangle (20,4.5); \node[anchor=east] at (18.9,4.3) {bitcoin reference rate window 19:00--20:00};
\end{tikzpicture}

\omcaption{Fixed moments of a northern-summer weekday in UTC, from the chapters of this book: continuous crypto trading with eight-hourly funding, the crypto
options expiry, the European day-ahead auction (12:00 CET), the LME's official-price publication (12:20--13:25 London), the US equity core session (9:30--16:00
New York) and the window of the New York bitcoin reference rate (15:00--16:00 New York). Schematic.}\label{fig:m3:a-map-of-all-markets:day}
\end{omfigure}

\begin{proposition}[Shifts to cover a round-the-clock book]\label{prop:m3:a-map-of-all-markets:shifts}
To staff a book that trades around the clock with shifts of $L$ hours that overlap by $o$ hours for handover, a desk needs $\lceil 24/(L - o)\rceil$ shifts a day;
the number of traders on each shift is the largest number of markets open at the same time during it.
\end{proposition}

\begin{proof}
Successive shifts start $L - o$ hours apart, so $n$ shifts cover $n(L - o)$ hours of the cycle, and $n$ must reach 24. A trader per open market is needed at the
busiest minute of each shift.
\end{proof}

\begin{omfigure}
\begin{tikzpicture}
\begin{axis}[width=11cm,height=4.6cm,ybar=0pt,bar width=4pt,
  xlabel={hour (UTC)},ylabel={desks open},
  tick label style={font=\small},label style={font=\small},
  xmin=-0.8,xmax=23.8,ymin=0,ymax=3.5,ytick={0,1,2,3},xtick={0,3,6,9,12,15,18,21},grid=major,grid style={gray!25},
  legend columns=2,legend style={font=\footnotesize,at={(0.5,-0.34)},anchor=north,draw=none,fill=none,/tikz/every even column/.append style={column sep=8pt}},
  table/col sep=comma]
\addplot[fill=omThm!50,draw=omThm] table[x=hour,y=weekday]{figdata/markets-3/29-a-map-of-all-markets/staffing.csv};\addlegendentry{weekday}
\addplot[fill=omDef!40,draw=omDef] table[x=hour,y=weekend]{figdata/markets-3/29-a-map-of-all-markets/staffing.csv};\addlegendentry{weekend}
\end{axis}
\end{tikzpicture}

\omcaption{Desks a firm must staff by hour, for the weekend problem's choices: crypto around the clock, European power 05:00--17:00 UTC and US markets
12:00--21:00 UTC on weekdays. Three shifts of nine hours overlapping by one cover the cycle (\cref{prop:m3:a-map-of-all-markets:shifts}); the busiest
hours need three traders. Illustrative desk hours. Data: the chapter's tutorial.}\label{fig:m3:a-map-of-all-markets:staffing}
\end{omfigure}

\section{Participants}

The same few kinds of firm appear in every market of these books, in different proportions: dealers (banks and their successors), who provide balance sheet and
take the other side of clients; proprietary trading firms, which trade their own capital, mostly in the order books and increasingly in the dealer markets
(One Quant Book 1); asset managers and hedge funds, who take views and pay for liquidity; physical players, the producers, consumers and trading houses of the
commodity chapters, whose trading follows their assets; and, in crypto, retail traders in large numbers, \omterm{def:m3:blockchains-for-traders:key}{token} issuers and protocols. Who is on the
other side decides a strategy's edge: in markets where hedgers and retail trade for reasons other than price, a market maker earns the spread; in markets of
professionals only, edge comes from information, speed or balance sheet.

\section{Dominant strategy types}

Across the markets of these books the strategies that pay fall into a few families, which the later books take up one by one. Market making and liquidity
provision: equities, futures, options, FX, crypto venues and pools (One Quant Books 10 and 11 on microstructure and high-frequency trading). Relative value and
arbitrage (One Quant Book 9): basis trades, \omterm{def:m3:spot-markets:cross}{cross-venue arbitrage}, spreads between related contracts, from Treasury futures to the \omterm{def:m3:spot-markets:kimchi}{kimchi premium}. Carry and risk premia (One Quant Books 8 and 9): funding
trades, \omterm{def:m3:forward-curves-storage-and-convenience-yield:roll}{roll yield}, selling volatility, the \omterm{def:m3:prediction-and-betting-markets:flb}{favourite--longshot bias}. Directional and systematic strategies (One Quant Book 8): trend and momentum in futures, value and
factor strategies in equities. And physical optimisation: storage, transport and generation in commodities, where the trade is an asset's operation. A firm's map is
the intersection of these families with the markets whose structure, hours and participants it can serve.

\section{Tutorial: the map and the desk}\label{tut:m3:a-map-of-all-markets:map}

\begin{tutorial}
\textbf{Goal.} Put market sizes on one logarithmic scale, compute turnover velocities, and plan the shifts that cover a firm's chosen markets around the clock.
\textbf{End state:} \cref{fig:m3:a-map-of-all-markets:sizes} and the numbers of the weekend problem.
\begin{tutsteps}
\item \textbf{The planner.} Fewest overlapping shifts covering every staffed minute, and traders per shift.
\omcode{code/firm/marketmap/firm_marketmap.py}{40}{61}{Shift planning over a day of market hours.}\label{lst:m3:a-map-of-all-markets:shifts}
\item \textbf{The desk.} \texttt{desks()} sets the firm's coverage choices; \texttt{shifts(9, 1)} plans a weekday and a weekend day.
\item \textbf{Run} \texttt{turnover\_velocity} for Treasuries and a perpetual, and \texttt{fig\_map.py}.
\end{tutsteps}
\textbf{What to change next.} Add an Asian desk and see whether it saves a shift; plan around daylight-saving changes; count trader-hours a week and compare with
the cost of an automated overnight shift with a human on call.
\end{tutorial}

\section{Build: the market registry}\label{bld:m3:a-map-of-all-markets:marketmap}

\begin{build}
\bfield{Purpose} The miniature firm must know, for every market it trades, when it is open, how it is structured and cleared, and how large it is; its operations
must plan staff and systems around those hours.
\bfield{Interface} \texttt{Market(name, daily\_usd\_bn, structure, cleared, hours)}; \texttt{every\_day(start, end, days)}; \texttt{is\_open},
\texttt{open\_markets(markets, weekday, minute)}; \texttt{staffed\_minutes}; \texttt{plan\_shifts(markets, weekday, length\_h, overlap\_h)}.
\bfield{Rules} Hours in UTC minutes per weekday, converted from local rules with daylight-saving time and holidays (from \texttt{firm.calendar} in a full
implementation); sizes stored with their source's measure; a round-the-clock market planned as a cycle.
\bfield{Acceptance tests} \texttt{code/firm/marketmap/tests/}: open queries on weekdays and weekends; three nine-hour shifts cover the cycle; one shift covers the
US session; no shift on a day with nothing open.
\bfield{Stretch} Daylight-saving transitions; holiday calendars per market; hand-over checklists generated from open positions.
\end{build}

\begin{omsources}
\item BIS, Triennial Central Bank Survey 2025 (One Quant Book 2's ledger); SIFMA, US Treasury and US fixed income securities statistics, September 2026;
FIA, ETD volume, December 2025 (One Quant Book 1's ledger).
\item Binance and Deribit public API endpoints, 24 September 2026.
\item NYSE, hours and calendars; this book's chapters 5, 8, 17 and 19 for the other fixed moments.
\item CME Group, trading hours page (FX Spot+ on Globex, 2025; Internet Archive copy); LME, data highlights 2025 (January 2026); EPEX SPOT, annual
trading results 2025 (19 January 2026).
\end{omsources}

\section{Exercises}

\begin{exercise}[$\star$]\label{exo:m3:a-map-of-all-markets:1}
What is the \omterm{def:m3:a-map-of-all-markets:velocity}{turnover velocity} of the US Treasury market from \cref{dat:m3:a-map-of-all-markets:sizes}?
\end{exercise}

\begin{exercise}[$\star$]\label{exo:m3:a-map-of-all-markets:2}
Convert to UTC in the northern summer: the US equity core session and the window of the New York bitcoin reference rate.
\end{exercise}

\begin{exercise}[$\star$]\label{exo:m3:a-map-of-all-markets:3}
Place three markets of your choice between the order-book pole and the dealer pole, and say why.
\end{exercise}

\begin{exercise}[$\star\star$]\label{exo:m3:a-map-of-all-markets:4}
How many shifts of eight hours, overlapping by one, does a round-the-clock desk need? Of twelve hours?
\end{exercise}

\begin{exercise}[$\star\star$]\label{exo:m3:a-map-of-all-markets:5}
Why is a comparison of FX turnover with one crypto venue's volume only indicative?
\end{exercise}

\begin{exercise}[$\star\star$]\label{exo:m3:a-map-of-all-markets:6}
What does a velocity above two for a \omterm{def:m3:perpetual-futures:perp}{perpetual future} say about who trades it?
\end{exercise}

\begin{exercise}[$\star\star\star$]\label{exo:m3:a-map-of-all-markets:7}
\emph{Coding.} With \texttt{shifts}, plan the tutorial's desks with nine-hour shifts on a weekday and on a Sunday. How many trader-shifts does each need?
\end{exercise}

\begin{exercise}[$\star\star\star$]\label{exo:m3:a-map-of-all-markets:8}
\emph{Find the flaw.} ``The biggest markets are the best opportunities.''
\end{exercise}

\section{Problem: The Round-the-Clock Desk}

\begin{problem}[{Weekend problem --- staffing crypto, European power and US markets}]\label{pb:m3:a-map-of-all-markets:1}
A firm trades crypto around the clock, European power from 05:00 to 17:00 UTC and US rates and equities from 12:00 to 21:00 UTC on weekdays. Shifts last nine
hours and overlap by one hour for handover; each market open needs one trader.

\medskip\noindent\textbf{Part I --- The shifts.}
\begin{enumerate}
\item How many shifts a day does the crypto book alone require?
\item When do the shifts start and end if the first starts at midnight UTC?
\item How many traders does each weekday shift need?
\item And at weekends?
\item How many trader-shifts a week is that?
\end{enumerate}

\medskip\noindent\textbf{Part II --- The moments.}
\begin{enumerate}[resume]
\item Which fixed moments of \cref{fig:m3:a-map-of-all-markets:day} fall in each shift?
\item Which shift handles the European day-ahead auction?
\item Which the US equity close and the bitcoin reference rate?
\item What should each handover cover?
\item How does daylight-saving time change the plan in winter?
\end{enumerate}

\medskip\noindent\textbf{Part III --- The alternatives.}
\begin{enumerate}[resume]
\item What would twelve-hour shifts save and cost?
\item Could an office in Asia replace the night shift?
\item What can be automated at night, and what must a human watch?
\item How would you cover the weekend?
\item What does the plan cost in salaries if a trader costs a stated amount a year?
\end{enumerate}

\medskip\noindent\textbf{Part IV --- Judgement.}
\begin{enumerate}[resume]
\item Which of the three markets should the firm drop if it had to, and why?
\item What does the map say about where to expand next?
\item How does \omterm{def:m3:a-map-of-all-markets:velocity}{turnover velocity} inform the choice?
\item State the \emph{named result}: the number of shifts and the traders per shift that cover the chosen markets with a one-hour overlap.
\item In one sentence: what is the map for?
\end{enumerate}
\end{problem}

\section{Interview questions}

\begin{interviewq}[$\star$ \iqroles{trader}]\label{iq:m3:a-map-of-all-markets:1}
Rank FX, US Treasuries, equities and crypto by daily turnover, and say what the ranking hides.
\end{interviewq}

\begin{interviewq}[$\star$ \iqroles{risk}]\label{iq:m3:a-map-of-all-markets:2}
What are the risks of a book handed between time zones?
\end{interviewq}

\begin{interviewq}[$\star\star$ \iqroles{researcher}]\label{iq:m3:a-map-of-all-markets:3}
How would you estimate the addressable profit pool of a new market for a market-making firm?
\end{interviewq}

\begin{interviewq}[$\star\star$ \iqroles{trader}]\label{iq:m3:a-map-of-all-markets:4}
Which market in these books would you enter with a small team and little capital, and why?
\end{interviewq}

\begin{interviewq}[$\star\star$ \iqroles{developer}]\label{iq:m3:a-map-of-all-markets:5}
Design a service that tells every system in a firm which markets are open, with holidays and daylight-saving time.
\end{interviewq}

\begin{interviewq}[$\star\star\star$ \iqroles{researcher, trader}]\label{iq:m3:a-map-of-all-markets:6}
Choose three markets from these books whose combination gives a firm the best use of the same technology, people and capital.
\end{interviewq}
